\documentclass[11pt]{article}
\usepackage[a4paper,margin=25mm]{geometry}
\usepackage{amsmath,amssymb,amsthm,hyperref}
\newtheorem{theorem}{Theorem}
\newtheorem{lemma}{Lemma}
\newtheorem{corollary}{Corollary}
\newcommand{\PP}{\mathcal P}
\title{Polynomial conditioning for a fixed number of separated Legendre bands}
\author{Proof independently reviewed by root, size\_bounds, and independent\_directions}
\date{30 September 2026}
\begin{document}\maketitle
All norms in this note are for the uniform probability measure on $[0,1]$.
Put
\[
 H_d(t)=\int_0^t P_d(2u-1)\,du\qquad(d\ge1).
\]
Thus $H_d'=P_d(2t-1)$ and $H_d(0)=H_d(1)=0$.
This note supplies conditioning for a possible fixed-number-of-small-dice
construction. It does not by itself construct that family.

\begin{lemma}\label{factor}
Let $D\ge d+1$, $p,q$ be real polynomials, $\deg p<d$ and
$\deg(p+H_dq)\le D$. Then
\[
 \|H_dq\|_2\ge c(D+1)^{-4}\|q\|_2,
 \qquad
 (\|p\|_2^2+\|q\|_2^2)^{1/2}
       \le C(D+1)^5\|p+H_dq\|_2.
\]
Consequently replacing $H_d$ by $\eta H_d$, $0<\eta\le1$, costs
at most an additional factor $\eta^{-1}$ in the second inequality.
\end{lemma}
\begin{proof}
The identity
\[
 H_d(t)=\frac{(x^2-1)P_d'(x)}{2d(d+1)},\qquad x=2t-1,
\]
shows that its $d+1$ simple zeros are the Gauss--Lobatto nodes,
including both endpoints. Their positive normalized quadrature weights
are
\[
 \omega_i=\frac1{d(d+1)P_d(2t_i-1)^2}.
\]
They sum to one, and the rule is exact through degree $2d-1$.
These are the classical Legendre--Lobatto formulas; they also follow
by integrating the cardinal polynomials and using Legendre orthogonality.
In particular,
\[
 |H_d'(t_i)|\ge[d(d+1)]^{-1/2}\ge c/d,
 \qquad \|H_d''\|_\infty\le C d^2.
\]
On an interval of radius $c d^{-3}$ around each zero, its derivative
therefore has magnitude at least $c/d$. These intervals can be made
disjoint, by Rolle's theorem and the same derivative bound.
There is exactly one critical point between consecutive zeros, and
no others. Thus $|H_d|$ is unimodal on each such interval.
For $a=c_0(D+1)^{-4}$, with sufficiently small absolute $c_0$,
\[
 |\{t:|H_d(t)|<a\}|\le C a d^2.
\]
The near-zero neighborhoods cover the whole sublevel set because their
outer boundary values are at least $c d^{-4}$, and $D\ge d$.
The endpoint neighborhoods are one-sided and satisfy the same estimate.

For a polynomial $q$ of degree at most $D$, the orthonormal Legendre
expansion gives $\|q\|_\infty\le(D+1)\|q\|_2$.
Hence at most $C c_0$ of its squared norm lies in this sublevel set.
Taking $c_0$ small gives the first asserted inequality.

Write $f=p+H_dq$. Since $2\deg p\le2d-2$, Lobatto exactness gives
\[
 \|p\|_2^2=\sum_i\omega_i p(t_i)^2
          =\sum_i\omega_i f(t_i)^2
          \le\|f\|_\infty^2\le(D+1)^2\|f\|_2^2.
\]
Thus $\|H_dq\|_2\le(D+2)\|f\|_2$. The multiplication lower bound
now proves the second inequality. Apply it to $\eta q$ for the
scaled assertion. Degree cancellation cannot hide a larger degree of
$q$, since $\deg p<d<\deg H_d$.
\end{proof}

\begin{theorem}[Lacunary product conditioning]\label{main}
Fix an integer $k\ge1$ and $r\ge0$. Choose integers $d_1,\ldots,d_k$
so that
\[
 d_i>r+\sum_{j<i}(d_j+1),
 \qquad D=r+\sum_{i=1}^k(d_i+1).
\]
For $0<\eta_i\le1$ and arbitrary polynomials $p_S\in\PP_r$,
indexed by subsets $S\subseteq[k]$, one has
\[
 \left(\sum_{S\subseteq[k]}\|p_S\|_2^2\right)^{1/2}
 \le C_k(D+1)^{5k}\left(\prod_{i=1}^k\eta_i^{-1}\right)
 \left\|\sum_{S\subseteq[k]}p_S(t)
                  \prod_{i\in S}\eta_iH_{d_i}(t)\right\|_2.
 \tag{1}
\]
In particular this representation is unique.
Taking $d_i=1+r+\sum_{j<i}(d_j+1)$ gives
$D=2^k(r+2)-2$. For fixed $k$, polynomially small $\eta_i$ therefore
produce a polynomial lower bound for the smallest singular value.
\end{theorem}
\begin{proof}
Decompose the displayed sum as $F_0+\eta_kH_{d_k}F_1$.
Both $F_0$ and $F_1$ have degree at most
$r+\sum_{j<k}(d_j+1)<d_k$. Lemma~\ref{factor} bounds the sum of
their squared norms by the squared norm of the full sum times
$C(D+1)^{10}\eta_k^{-2}$. Apply the same argument recursively to
each branch. At most $k$ levels occur, all degrees are at most $D$,
and each level contributes its corresponding $\eta_i^{-1}$.
This proves (1), including the harmless constant depending only on $k$.
\end{proof}

\begin{corollary}[Products of monotone ranks]\label{ranks}
Put $R_i(t)=t-\eta_iH_{d_i}(t)$, with $0<\eta_i\le1/2$.
These maps are strictly increasing from $0$ to $1$.
Suppose
\[
 d_i>r+k+\sum_{j<i}(d_j+1),\qquad
 D=r+k+\sum_i(d_i+1).
\]
Then the family of functions
\[
 \{\phi_a(t)\prod_{i\in S}R_i(t):0\le a\le r,\ S\subseteq[k]\}
\]
is linearly independent, and its least singular value in $L^2[0,1]$
is at least
$c_k(D+1)^{-5k}\prod_i\eta_i$.
Here $\phi_a$ are the ordinary orthonormal shifted Legendre polynomials.
\end{corollary}
\begin{proof}
Since $R_i'=1-\eta_iP_{d_i}(2t-1)\ge1/2$, monotonicity is immediate.
For $F=\sum_S p_S\prod_{i\in S}R_i$, expansion gives
\[
 F=\sum_T(-1)^{|T|}\left(\prod_{i\in T}\eta_iH_{d_i}\right)q_T,
 \qquad q_T=\sum_{S\supseteq T}t^{|S\setminus T|}p_S.
\]
Each $q_T$ has degree at most $r+k$. The inverse triangular identity is
\[
 p_S=\sum_{T\supseteq S}(-t)^{|T\setminus S|}q_T.
\]
As multiplication by a power of $t$ is an $L^2$ contraction, both
transformations have norm bounded by a constant depending only on $k$.
Theorem~\ref{main}, applied with degree budget $r+k$, proves the claim.
\end{proof}

\begin{corollary}[Discrete sampling]
For the spaces in Theorem~\ref{main} or Corollary~\ref{ranks}, the same
conditioning estimates, up to an absolute factor, hold in empirical
$L^2$ on any node set with star discrepancy at most $c(D+1)^{-2}$.
In particular they hold on each of a fixed number of interlaced groups
of a sufficiently fine perturbed midpoint grid.
\end{corollary}
\begin{proof}
Every function $F$ involved is a polynomial of degree at most $D$.
The discrepancy inequality and the $L^2$ Markov inequality give
\[
 \left|\frac1N\sum_qF(t_q)^2-\int_0^1F^2\right|
 \le\operatorname{disc}(t)\int_0^1|(F^2)'|
 \le C(D+1)^2\operatorname{disc}(t)\|F\|_2^2.
\]
For sufficiently small $c$ this is at most half the squared norm.
One elementary proof of the needed Markov inequality expands $F$ in
Legendre polynomials and uses
$\|\phi_j'\|_2^2=2(2j+1)j(j+1)$, followed by Cauchy--Schwarz.
\end{proof}
\end{document}
